\documentclass[letterpaper,11pt]{article}
\usepackage[top=0.8in, bottom=0.8in, left=0.9in, right=0.9in]{geometry}
\usepackage[hidelinks]{hyperref}
\usepackage{lmodern}
\usepackage{parskip}
\usepackage{microtype}
\setlength{\parskip}{6pt}
\setlength{\parindent}{0pt}
\begin{document}
\begin{flushleft}
\textbf{\large Ajay Venkatesh} \\
Brooklyn, NY \\
+1 516 585 9013 \\
\href{mailto:av3855@nyu.edu}{av3855@nyu.edu}
\end{flushleft}
\vspace{10pt}
May 21, 2026
\vspace{6pt}

Hiring Manager \\
SpaceX

\vspace{10pt}

Dear Hiring Manager,

My name is Ajay Venkatesh. I'm finishing my M.S. in Computer Engineering at NYU, with production software experience at PwC in Bengaluru, a summer SDE internship at Amazon, and ongoing on-campus work at NYU's Novel AI Technologies. I'm writing about the ML Surrogate Modeling Engineer role on the Vehicle Engineering team.

The role's combination, developing AI surrogate models that accelerate engineering simulation workflows for launch vehicles and spacecraft, sits at the intersection of where my graduate coursework and applied AI work have been pushing. I've worked across \textbf{PyTorch}, training and evaluating models in production contexts. At Amazon I deployed a \textbf{Retrieval-Augmented Generation (RAG)} pipeline with an \textbf{agentic workflow} that reasons over operational logs to generate context-aware summaries, and shipped a \textbf{Model Context Protocol (MCP)} server that materially reduced query time and compute costs. Deploying production-grade AI models with attention to inference latency and reliability is the shape of work I want to keep doing.

On modeling and scientific computing: my graduate coursework includes \textbf{Machine Learning}, \textbf{Deep Learning}, and \textbf{Big Data}, plus a \textbf{PySpark / HDFS} project over millions of flight records that trained Random Forest and Gradient Boosted Tree models with rigorous EDA. I've also worked on RoBERTa fine-tuning with custom LoRA implementations and ResNet image classifiers under strict parameter budgets, training stability practices that carry over to surrogate-model development.

On scalable data pipelines: at NYU I engineered \textbf{ETL pipelines} transforming \textbf{NoSQL} into \textbf{PostgreSQL} with concurrency control, and provisioned scalable \textbf{AWS} infrastructure for data preprocessing. Building reliable data pipelines that feed production models is the kind of plumbing surrogate-modeling workflows depend on.

Stack-wise: \textbf{Python}, \textbf{PyTorch}, \textbf{NumPy}, \textbf{Pandas}, \textbf{Linux}, plus statistics fundamentals from coursework. Comfortable with software engineering best practices, code review, and writing production-ready ML code.

What pulls me toward SpaceX Vehicle Engineering is the stakes. Engineering for systems where mistakes happen in real units, on real hardware, with consequences that translate into missions and lives is exactly the kind of work that demands the most careful engineering. The chance to build models that actually accelerate the development of rockets and spacecraft is rare.

I'd welcome the chance to discuss further. Thanks for considering my application.

\vspace{12pt}
Sincerely, \\
Ajay Venkatesh

\end{document}